% part: intuitionistic-logic 在直觉主义逻辑中，真是依据“构造性证明”或“证实”来解释的：一个命题为真，当且仅当它（已经或能够）被证明。我们如何确认自己证明了一个命题呢？我们可能有一个直接证明，或者一个通过可靠方法得到的证明。对于复合陈述，其各组成部分的证明必须以适当的方式连接起来，才能构成该复合陈述的证明。这构成了与经典的真值条件语义学的根本背离。在经典语义学中，复合陈述的真值是各组成部分真值的函数，而真值本身并不通过证明或证实来解释。% chapter: soundness-completeness 直觉主义逻辑的 Kripke 语义是使这种证明条件思想在数学上精确化的一种方式。它由 % section: truth-lemma Brouwer 等人在 20 世纪早期发展起来。

\documentclass[../../../include/open-logic-section]{subfiles}

\begin{document}

\olfileid{int}{sc}{tru}

\olsection{真值引理}

下述引理将典范模型中的满足性与哪些公式是素集 $\Delta$ 的元素联系了起来。

\begin{lem}\ollabel{lem:truth}
  如果 $\Delta$ 是素的，那么 $\mSat{M(\Delta)}{!A}[\sigma]$ 当且仅当
  $\Delta(\sigma) \Proves !A$。
\end{lem}

\begin{proof}
  对~$!A$ 进行归纳。
  \begin{enumerate}
    \item \indcase{!A}{\lfalse}{由于 $\Delta(\sigma)$ 是素的，它
      是一致的，所以 $\Delta(\sigma) \Proves/ \indfrm$。根据
      定义，$\mSat/{M(\Delta)}{\indfrm}[\sigma]$。}
      \item \indcase{!A}{p}{根据 $\mSat{{}}{}$ 的定义，
        $\mSat{M(\Delta)}{\indfrm}[\sigma]$ 当且仅当 $\sigma \in V(p)$，
        即 $\Delta(\sigma) \Proves \indfrm$。}
      \item \indcase!{!A}{\lnot !B}{}
      \item \indcase{!A}{!B \land
        !C}{$\mSat{M(\Delta)}{\indfrm}[\sigma]$ 当且仅当
        $\mSat{M(\Delta)}{!B}[\sigma]$ 且
        $\mSat{M(\Delta)}{!C}[\sigma]$。由归纳假设，
        $\mSat{M(\Delta)}{!B}[\sigma]$ 当且仅当 $\Delta(\sigma) \Proves
        !B$，对~$!C$ 同理。而 $\Delta(\sigma) \Proves !B$
        且 $\Delta(\sigma) \Proves !C$ 当且仅当 $\Delta(\sigma)
        \Proves \indfrm$。}
      \item \indcase{!A}{!B \lor
        !C}{$\mSat{M(\Delta)}{\indfrm}[\sigma]$ 当且仅当
        $\mSat{M(\Delta)}{!B}[\sigma]$ 或
        $\mSat{M(\Delta)}{!C}[\sigma]$。由归纳假设，
        这等价于 $\Delta(\sigma) \Proves !B$ 或
        $\Delta(\sigma) \Proves !C$。我们需要证明这又等价于 $\Delta(\sigma) \Proves \indfrm$。
        从左到右的方向是显然的。从右到左的方向则由 $\Delta(\sigma)$ 是素集可得。}
      \item \indcase{!A}{!B \lif !C}{先证从左到右方向的逆否命题：假设 $\Delta(\sigma) \Proves/ !B
        \lif !C$。
        那么也有 $\Delta(\sigma) \cup \{!B\} \Proves/
        !C$。由于 $\tuple{!B, !C}$ 对
        某个~$n$ 是 $\tuple{!B_n, !C_n}$，我们有 $\Delta(\sigma.n) = (\Delta(\sigma) \cup
        \{!B\})^*$，且 $\Delta(\sigma.n) \Proves !B$ 但 $\Delta(\sigma.n) \Proves/
        !C$。由归纳假设，$\mSat{M(\Delta)}{!B}[\sigma.n]$
        且 $\mSat/{M(\Delta)}{!C}[\sigma.n]$。由于
        $R\sigma(\sigma.n)$，这意味着
        $\mSat/{M(\Delta)}{\indfrm}[\sigma]$。

        现在假设 $\Delta(\sigma) \Proves !B \lif !C$，并设
        $R\sigma\sigma'$。由于 $\Delta(\sigma) \subseteq
        \Delta(\sigma')$，我们有：若 $\Delta(\sigma') \Proves
        !B$，则 $\Delta(\sigma') \Proves !C$。换言之，对每个
        满足 $R\sigma\sigma'$ 的 $\sigma'$，要么
        $\Delta(\sigma') \Proves/ !B$，要么 $\Delta(\sigma') \Proves
        !C$。由归纳假设，这意味着每当
        $R\sigma\sigma'$，要么 $\mSat/{M(\Delta)}{!B}[\sigma']$，要么
        $\mSat{M(\Delta)}{!C}[\sigma']$，即
        $\mSat{M(\Delta)}{\indfrm}[\sigma]$。}
  \end{enumerate}
\end{proof}

\end{document}